% Options for packages loaded elsewhere
\PassOptionsToPackage{unicode}{hyperref}
\PassOptionsToPackage{hyphens}{url}
\documentclass[
  polish,
  11pt,
  a4paper,
]{article}
\usepackage{xcolor}
\usepackage[margin=2cm]{geometry}
\usepackage{amsmath,amssymb}
\setcounter{secnumdepth}{-\maxdimen} % remove section numbering
\usepackage{iftex}
\ifPDFTeX
  \usepackage[T1]{fontenc}
  \usepackage[utf8]{inputenc}
  \usepackage{textcomp} % provide euro and other symbols
\else % if luatex or xetex
  \usepackage{unicode-math} % this also loads fontspec
  \defaultfontfeatures{Scale=MatchLowercase}
  \defaultfontfeatures[\rmfamily]{Ligatures=TeX,Scale=1}
\fi
\usepackage{lmodern}
\ifPDFTeX\else
  % xetex/luatex font selection
\fi
% Use upquote if available, for straight quotes in verbatim environments
\IfFileExists{upquote.sty}{\usepackage{upquote}}{}
\IfFileExists{microtype.sty}{% use microtype if available
  \usepackage[]{microtype}
  \UseMicrotypeSet[protrusion]{basicmath} % disable protrusion for tt fonts
}{}
\makeatletter
\@ifundefined{KOMAClassName}{% if non-KOMA class
  \IfFileExists{parskip.sty}{%
    \usepackage{parskip}
  }{% else
    \setlength{\parindent}{0pt}
    \setlength{\parskip}{6pt plus 2pt minus 1pt}}
}{% if KOMA class
  \KOMAoptions{parskip=half}}
\makeatother
\usepackage{color}
\usepackage{fancyvrb}
\newcommand{\VerbBar}{|}
\newcommand{\VERB}{\Verb[commandchars=\\\{\}]}
\DefineVerbatimEnvironment{Highlighting}{Verbatim}{commandchars=\\\{\}}
% Add ',fontsize=\small' for more characters per line
\usepackage{framed}
\definecolor{shadecolor}{RGB}{248,248,248}
\newenvironment{Shaded}{\begin{snugshade}}{\end{snugshade}}
\newcommand{\AlertTok}[1]{\textcolor[rgb]{0.94,0.16,0.16}{#1}}
\newcommand{\AnnotationTok}[1]{\textcolor[rgb]{0.56,0.35,0.01}{\textbf{\textit{#1}}}}
\newcommand{\AttributeTok}[1]{\textcolor[rgb]{0.13,0.29,0.53}{#1}}
\newcommand{\BaseNTok}[1]{\textcolor[rgb]{0.00,0.00,0.81}{#1}}
\newcommand{\BuiltInTok}[1]{#1}
\newcommand{\CharTok}[1]{\textcolor[rgb]{0.31,0.60,0.02}{#1}}
\newcommand{\CommentTok}[1]{\textcolor[rgb]{0.56,0.35,0.01}{\textit{#1}}}
\newcommand{\CommentVarTok}[1]{\textcolor[rgb]{0.56,0.35,0.01}{\textbf{\textit{#1}}}}
\newcommand{\ConstantTok}[1]{\textcolor[rgb]{0.56,0.35,0.01}{#1}}
\newcommand{\ControlFlowTok}[1]{\textcolor[rgb]{0.13,0.29,0.53}{\textbf{#1}}}
\newcommand{\DataTypeTok}[1]{\textcolor[rgb]{0.13,0.29,0.53}{#1}}
\newcommand{\DecValTok}[1]{\textcolor[rgb]{0.00,0.00,0.81}{#1}}
\newcommand{\DocumentationTok}[1]{\textcolor[rgb]{0.56,0.35,0.01}{\textbf{\textit{#1}}}}
\newcommand{\ErrorTok}[1]{\textcolor[rgb]{0.64,0.00,0.00}{\textbf{#1}}}
\newcommand{\ExtensionTok}[1]{#1}
\newcommand{\FloatTok}[1]{\textcolor[rgb]{0.00,0.00,0.81}{#1}}
\newcommand{\FunctionTok}[1]{\textcolor[rgb]{0.13,0.29,0.53}{\textbf{#1}}}
\newcommand{\ImportTok}[1]{#1}
\newcommand{\InformationTok}[1]{\textcolor[rgb]{0.56,0.35,0.01}{\textbf{\textit{#1}}}}
\newcommand{\KeywordTok}[1]{\textcolor[rgb]{0.13,0.29,0.53}{\textbf{#1}}}
\newcommand{\NormalTok}[1]{#1}
\newcommand{\OperatorTok}[1]{\textcolor[rgb]{0.81,0.36,0.00}{\textbf{#1}}}
\newcommand{\OtherTok}[1]{\textcolor[rgb]{0.56,0.35,0.01}{#1}}
\newcommand{\PreprocessorTok}[1]{\textcolor[rgb]{0.56,0.35,0.01}{\textit{#1}}}
\newcommand{\RegionMarkerTok}[1]{#1}
\newcommand{\SpecialCharTok}[1]{\textcolor[rgb]{0.81,0.36,0.00}{\textbf{#1}}}
\newcommand{\SpecialStringTok}[1]{\textcolor[rgb]{0.31,0.60,0.02}{#1}}
\newcommand{\StringTok}[1]{\textcolor[rgb]{0.31,0.60,0.02}{#1}}
\newcommand{\VariableTok}[1]{\textcolor[rgb]{0.00,0.00,0.00}{#1}}
\newcommand{\VerbatimStringTok}[1]{\textcolor[rgb]{0.31,0.60,0.02}{#1}}
\newcommand{\WarningTok}[1]{\textcolor[rgb]{0.56,0.35,0.01}{\textbf{\textit{#1}}}}
\usepackage{graphicx}
\makeatletter
\newsavebox\pandoc@box
\newcommand*\pandocbounded[1]{% scales image to fit in text height/width
  \sbox\pandoc@box{#1}%
  \Gscale@div\@tempa{\textheight}{\dimexpr\ht\pandoc@box+\dp\pandoc@box\relax}%
  \Gscale@div\@tempb{\linewidth}{\wd\pandoc@box}%
  \ifdim\@tempb\p@<\@tempa\p@\let\@tempa\@tempb\fi% select the smaller of both
  \ifdim\@tempa\p@<\p@\scalebox{\@tempa}{\usebox\pandoc@box}%
  \else\usebox{\pandoc@box}%
  \fi%
}
% Set default figure placement to htbp
\def\fps@figure{htbp}
\makeatother
\ifLuaTeX
\usepackage[bidi=basic,shorthands=off]{babel}
\else
\usepackage[bidi=default,shorthands=off]{babel}
\fi
\ifLuaTeX
  \usepackage{selnolig} % disable illegal ligatures
\fi
\setlength{\emergencystretch}{3em} % prevent overfull lines
\providecommand{\tightlist}{%
  \setlength{\itemsep}{0pt}\setlength{\parskip}{0pt}}
\usepackage{fontspec}
\defaultfontfeatures{Ligatures=TeX}
\usepackage{amsmath,amssymb}
\usepackage{booktabs}
\usepackage{array}
\usepackage{geometry}
\usepackage{enumitem}
\usepackage{xcolor}
\usepackage{tcolorbox}
\usepackage{fancyhdr}
\usepackage{titlesec}
\usepackage{multicol}

\geometry{margin=1.8cm, top=2.5cm, bottom=2cm}
\pagestyle{fancy}
\fancyhf{}
\rhead{\small Projektowanie badań — 2026/27}
\lhead{\small Zarządzanie informacją | ISI UJ}
\cfoot{\thepage}

\titleformat{\section}{\large\bfseries\color{blue!70!black}}{\thesection.}{0.5em}{}
\titleformat{\subsection}{\normalsize\bfseries}{\thesubsection.}{0.5em}{}
\titleformat{\subsubsection}{\small\bfseries\color{gray!70!black}}{\thesubsubsection.}{0.5em}{}
\titlespacing*{\section}{0pt}{1.5ex}{1ex}
\titlespacing*{\subsection}{0pt}{1ex}{0.5ex}

\setlist[itemize]{nosep, leftmargin=1.5em}
\setlist[enumerate]{nosep, leftmargin=1.5em}

\tcbset{boxrule=0.5pt, left=4pt, right=4pt, top=4pt, bottom=4pt, fonttitle=\bfseries\small}

\newtcolorbox{definicja}[1][]{colback=blue!5, colframe=blue!50!black, title={#1}}
\newtcolorbox{uwagabox}[1][]{colback=orange!10, colframe=orange!70!black, title={#1}}
\newtcolorbox{wzorbox}[1][]{colback=gray!5, colframe=gray!50!black, title={#1}}
\newtcolorbox{przykladbox}[1][]{colback=purple!5, colframe=purple!50!black, title={#1}}

\usepackage{bookmark}
\IfFileExists{xurl.sty}{\usepackage{xurl}}{} % add URL line breaks if available
\urlstyle{same}
% fallback for those not using the hyperref driver hyperxmp:
\makeatletter
\@ifundefined{xmpquote}{\newcommand{\xmpquote}[1]{#1}}{}
\makeatother
\hypersetup{
  pdftitle={Handout C05 --- Niepewność średniej{,} bootstrap i intuicja permutacji},
  pdfauthor={Marek Deja},
  pdflang={pl-PL},
  hidelinks,
  pdfcreator={LaTeX via pandoc}}

\title{Handout C05 --- Niepewność średniej, bootstrap i intuicja
permutacji}
\usepackage{etoolbox}
\makeatletter
\providecommand{\subtitle}[1]{% add subtitle to \maketitle
  \apptocmd{\@title}{\par {\large #1 \par}}{}{}
}
\makeatother
\subtitle{Zarządzanie informacją --- część ilościowa --- 2026/27}
\author{Marek Deja}
\date{2026/27}

\begin{document}
\maketitle

\section{C05 --- 90 minut}\label{c05-90-minut}

10 min odtworzenie → 20 min estymacja/CI → 20 min bootstrap → 10 min
permutacja → 20 min Z05 → 10 min odbiór. Własny indeks z reguł C02/C04,
N po wybraniu ważnych wyników.

\section{Estymata i CI}\label{estymata-i-ci}

\(s\) opisuje osoby, \(SE=s/\sqrt{n}\) niepewność średniej przy
niezależnych obserwacjach. CI t: \(\bar{x}\pm t_{0{,}975,n-1}SE\). Około
95\% przedziałów obejmowałoby parametr przy powtarzaniu właściwego
próbkowania; to nie zakres 95\% odpowiedzi.

Obliczamy indeks i zachowujemy tylko ważne wyniki. Poziom ufności
odczytujemy jawnie, a domyślny test względem zera pomijamy. Odczytaj N,
średnią oraz granice.

\begin{Shaded}
\begin{Highlighting}[]
\NormalTok{d }\OtherTok{\textless{}{-}}\NormalTok{ badaniaZI}\SpecialCharTok{::}\FunctionTok{przygotuj\_ankiete}\NormalTok{(badaniaZI}\SpecialCharTok{::}\FunctionTok{dane\_przykladowe}\NormalTok{())}\SpecialCharTok{$}\NormalTok{dane}
\NormalTok{p }\OtherTok{\textless{}{-}}\NormalTok{ d[, }\FunctionTok{paste0}\NormalTok{(}\StringTok{"pozycja\_"}\NormalTok{, }\DecValTok{1}\SpecialCharTok{:}\DecValTok{6}\NormalTok{)]}
\NormalTok{p}\SpecialCharTok{$}\NormalTok{pozycja\_3 }\OtherTok{\textless{}{-}}\NormalTok{ badaniaZI}\SpecialCharTok{::}\FunctionTok{odwroc\_pozycje}\NormalTok{(p}\SpecialCharTok{$}\NormalTok{pozycja\_3)}
\NormalTok{x }\OtherTok{\textless{}{-}}\NormalTok{ badaniaZI}\SpecialCharTok{::}\FunctionTok{indeks\_ankiety}\NormalTok{(p, }\AttributeTok{minimum =} \DecValTok{5}\NormalTok{)}
\NormalTok{x }\OtherTok{\textless{}{-}}\NormalTok{ x[}\SpecialCharTok{!}\FunctionTok{is.na}\NormalTok{(x)]}
\FunctionTok{c}\NormalTok{(}\AttributeTok{n =} \FunctionTok{length}\NormalTok{(x), }\AttributeTok{srednia =} \FunctionTok{mean}\NormalTok{(x), }\AttributeTok{se =}\NormalTok{ stats}\SpecialCharTok{::}\FunctionTok{sd}\NormalTok{(x) }\SpecialCharTok{/} \FunctionTok{sqrt}\NormalTok{(}\FunctionTok{length}\NormalTok{(x)))}
\end{Highlighting}
\end{Shaded}

\begin{verbatim}
##            n      srednia           se
## 147.00000000   3.29229025   0.05994961
\end{verbatim}

\begin{Shaded}
\begin{Highlighting}[]
\NormalTok{stats}\SpecialCharTok{::}\FunctionTok{t.test}\NormalTok{(x, }\AttributeTok{conf.level =}\NormalTok{ .}\DecValTok{95}\NormalTok{)}\SpecialCharTok{$}\NormalTok{conf.int}
\end{Highlighting}
\end{Shaded}

\begin{verbatim}
## [1] 3.173809 3.410771
## attr(,"conf.level")
## [1] 0.95
\end{verbatim}

\section{Bootstrap średniej}\label{bootstrap-ux15bredniej}

Losujemy respondentów ze zwracaniem, zachowując N. B oznacza repliki
komputera, nie nowe osoby. Kwantyle 2,5\% i 97,5\% dają proste
przybliżenie percentylowego CI.

\begin{Shaded}
\begin{Highlighting}[]
\FunctionTok{set.seed}\NormalTok{(}\DecValTok{2027}\NormalTok{)}
\NormalTok{B }\OtherTok{\textless{}{-}} \DecValTok{1999}
\NormalTok{boot }\OtherTok{\textless{}{-}} \FunctionTok{replicate}\NormalTok{(B, }\FunctionTok{mean}\NormalTok{(x[}\FunctionTok{sample.int}\NormalTok{(}\FunctionTok{length}\NormalTok{(x), }\FunctionTok{length}\NormalTok{(x), }\AttributeTok{replace =} \ConstantTok{TRUE}\NormalTok{)]))}
\NormalTok{stats}\SpecialCharTok{::}\FunctionTok{quantile}\NormalTok{(boot, }\FunctionTok{c}\NormalTok{(.}\DecValTok{025}\NormalTok{, .}\DecValTok{975}\NormalTok{))}
\end{Highlighting}
\end{Shaded}

\begin{verbatim}
##     2.5%    97.5%
## 3.173424 3.406587
\end{verbatim}

Rekrutacja, jakość pytań i zależność osób nadal wymagają oceny. Więcej
replik nie naprawia systematycznego błędu badania. Wąski CI oznacza
precyzję w modelu, a nie gwarancję trafnego pomiaru.

\section{Permutacja --- demonstracja}\label{permutacja-demonstracja}

Tasowanie etykiet grup przy zachowanych liczebnościach buduje rozkład
zerowy różnicy. Wymienialność wymaga uzasadnienia; brak normalności nie
zapewnia jej automatycznie. Dwustronne Monte Carlo:
\((1+\#\{|T_b|\geq|T_{obs}|\})/(B+1)\), więc skończona symulacja nie
daje p=0. Ta analiza nie jest osobnym wymaganiem Z05.

\section{Z05 --- dowody i rubryka}\label{z05-dowody-i-rubryka}

\texttt{analiza.R}, \texttt{odpowiedzi.md},
\texttt{wyniki/estymacja.csv}, \texttt{bootstrap.csv},
\texttt{bootstrap.png}. Dwie metody 95\%, N, średnia, SD/SE; B=1999,
ziarno 2027; histogram replik z N/B/scenariuszem i źródłem syntetycznym.
W MD: względne linki, znaczenie CI, B a N, precyzja, znaczenie
praktyczne i ograniczenie próby.

3 pkt estymata/CI/N + 2 pkt bootstrap/powtarzalność + 3 pkt własne
wyjaśnienie CI + 2 pkt precyzja/ograniczenie. Pełne poziomy:
\texttt{rubryka("Z05")}. Termin: początek C06. Zapisz, uruchom skrypt od
początku, sprawdź, oddaj i odczytaj potwierdzenie funkcjami R.

\section{Notatki}\label{notatki}

N: \ldots\ldots\ldots.. Średnia: \ldots\ldots\ldots.. SD:
\ldots\ldots\ldots.. SE: \ldots\ldots\ldots.. CI t:
\ldots\ldots\ldots\ldots\ldots\ldots\ldots.

CI bootstrap:
\ldots\ldots\ldots\ldots\ldots\ldots\ldots\ldots\ldots\ldots{} Dlaczego
B nie jest N?
\ldots\ldots\ldots\ldots\ldots\ldots\ldots\ldots\ldots\ldots\ldots..

Ograniczenie próby:
\ldots\ldots\ldots\ldots\ldots\ldots\ldots\ldots\ldots\ldots\ldots\ldots\ldots\ldots\ldots\ldots\ldots\ldots\ldots\ldots\ldots\ldots\ldots\ldots\ldots\ldots\ldots..

\end{document}
